\chapter{Come si modifica il metodo di cui sopra per calcolare la 
distanza di Hamming (se v e w hanno stessa lunghezza) e la 
longest common subsequence tra v e w}

Per adattare il metodo di calcolo della distanza di edit per ottenere la distanza 
di Hamming o Longest common subsequence (LCS), 
è necessario modificare le \textbf{operazioni consentite}, 
e di conseguenza la definizione della \textbf{funzione di costo} $\delta$ 
e l'algoritmo di \textbf{programmazione dinamica}.

\paragraph{Distanza di Hamming}

La \textbf{distanza di Hamming} si applica quando le due stringhe hanno la stessa 
lunghezza e si calcola confrontando i caratteri nelle stesse posizioni. 
Ogni posizione diversa contribuisce con costo $1$, mentre ogni posizione uguale 
contribuisce con costo $0$.

\begin{equation}
    d_H(v, w) = \text{numero di posizioni in cui } v \text{ e } w \text{ differiscono}
\end{equation}

\textbf{Esempio}

\begin{align*}
    S_1 &= ACCGT \\
    S_2 &= AGCGT
\end{align*}

Le due stringhe differiscono solo nella seconda posizione, dove $C \neq G$. Quindi:
\begin{equation}
    d_H(S_1,S_2)=1
\end{equation}

Se invece:
\begin{align*}
    S_1 &= ACCGT \\
    S_2 &= ACCGT
\end{align*}

non ci sono posizioni diverse, quindi:
\begin{equation}
    d_H(S_1,S_2)=0
\end{equation}

Per calcolare $d_H(v, w)$ utilizzando una griglia di edit definita per la distanza 
di edit, ma bisogna apportare alcune modifiche:
\begin{itemize}
    \item La griglia di edit sarà una matrice di dimensione $(n+1) \times (n+1)$, poiché le stringhe devono avere la stessa lunghezza.
    \item Gli archi consentiti saranno solo quelli diagonali, poiché non sono permesse operazioni di inserzione o delezione. 
    \item La funzione di costo $\delta$ sarà definita come:
    \begin{itemize}
        \item \textbf{match}: $\delta(v_i, w_j) = 0$ se $v_i = w_j$ 
        \item \textbf{mismatch}: $\delta(v_i, w_j) = 1$ se $v_i \neq w_j$
        \item \textbf{indel}: $\delta(v_i, \text{-})=\delta(\text{-}, w_j)=n$ se uno dei due simboli è un gap.
    \end{itemize}
\end{itemize}

Il \textbf{costo $n$} assegnato agli indel serve a penalizzare fortemente inserzioni e
delezioni. Infatti, poiché le due stringhe hanno la stessa lunghezza, il costo
massimo della distanza di Hamming è $n$, cioè il caso in cui tutti i caratteri
siano diversi. Se invece un allineamento usa un gap, per mantenere le due
sequenze della stessa lunghezza dovrà introdurre almeno due indel, con costo
totale almeno $2n$. Poiché $2n > n$, un allineamento con gap non potrà mai essere
ottimo. In questo modo l'algoritmo è forzato ad allineare solo caratteri in
\textbf{posizione corrispondente} (movimenti diagonali), e il costo finale coincide con il numero di mismatch.

Per calcolare il costo minimo si può usare ancora la programmazione dinamica,
come nel caso della distanza di edit. Indichiamo con $S(i,j)$ il costo minimo
per allineare il prefisso $v[1..i]$ con il prefisso $w[1..j]$.

I \textbf{casi base} diventano:
\begin{itemize}
    \item $S(0,0)=0$
    \item $S(i,0)=i \cdot n$ per $0 \leq i \leq n$
    \item $S(0,j)=j \cdot n$ per $0 \leq j \leq n$
\end{itemize}

Infatti, per allineare un prefisso non vuoto con una stringa vuota, è necessario
inserire solo gap, e ogni gap ha costo $n$.

Successivamente, si riempie il resto della griglia con la seguente
\textbf{equazione di ricorrenza}:
\begin{equation}
    S(i,j)=\min
    \begin{cases}
        S(i-1,j)+n & \text{(Delezione)} \\
        S(i,j-1)+n & \text{(Inserzione)} \\
        S(i-1,j-1)+\delta(v_i,w_j) & \text{(Match/Mismatch)}
    \end{cases}
\end{equation}

dove:

\begin{equation}
    \delta(v_i,w_j)=
    \begin{cases}
        0 & \text{se } v_i=w_j \\
        1 & \text{se } v_i \neq w_j
    \end{cases}
\end{equation}

Il valore nel nodo finale:
\begin{equation}
    S(n,n)
\end{equation}

coincide quindi con la distanza di Hamming:

\begin{equation}
     S(n,n)=d_H(v,w)
\end{equation}

La \textbf{matrice di allineamento} che si ottiene non contiene gap, ma solo colonne del tipo:

\[
\begin{pmatrix}
v_i \\
w_i
\end{pmatrix}
\]

per ogni posizione $i$. Ogni colonna ha costo $0$ se i due simboli sono uguali e
costo $1$ se sono diversi. La somma dei costi delle colonne è quindi esattamente
il numero di mismatch tra $v$ e $w$.
In questo modo, la distanza di Hamming può essere vista come un caso particolare
della distanza di edit in cui inserzioni e delezioni sono rese non convenienti
tramite un costo molto alto.

\paragraph{Longest Common Subsequence (LCS)}
Date due sequenze $v = v_1 v_2 \ldots v_n$ e $w = w_1 w_2 \ldots w_m$ 
la \textbf{longest common subsequence (LCS)} di $v$ e $w$ è la sequenza di posizioni
$1 \leq i_1 < i_2 < \ldots < i_k \leq n$ e $1 \leq j_1 < j_2 < \ldots < j_k \leq m$
tale che la $i$-esima posizione di $v$ e la $j$-esima posizione di $w$ contengono lo stesso simbolo, 
ovvero $v_{i_t} = w_{j_t}$ per ogni $t = 1, 2, \ldots, k$. 
La lunghezza della longest common subsequence è indicata con $LCS(v,w)$.

\textbf{Esempio}

Consideriamo:
\begin{align*}
    v &= ATCTGATC \\
    w &= TGCATAC
\end{align*}

Un possibile allineamento tra $v$ e $w$ è:
\[
A =
\begin{pmatrix}
A & \textcolor{red}{T} & \text{-} & \textcolor{red}{C} & \text{-} & \textcolor{red}{T} & G & \textcolor{red}{A} & T & \textcolor{red}{C} \\
\text{-} & \textcolor{red}{T} & G & \textcolor{red}{C} & A & \textcolor{red}{T} & \text{-} & \textcolor{red}{A} & \text{-} & \textcolor{red}{C}
\end{pmatrix}
\]

Le colonne evidenziate in rosso sono i \textbf{match}, cioè le colonne in cui i due
simboli allineati sono uguali. Il cammino corrispondente nella griglia è:
\[
\begin{aligned}
(0,0) &\rightarrow (1,0) \rightarrow \textcolor{red}{(2,1)}
\rightarrow (2,2) \rightarrow \textcolor{red}{(3,3)} \rightarrow (3,4) \\
&\rightarrow \textcolor{red}{(4,5)} \rightarrow (5,5)
\rightarrow \textcolor{red}{(6,6)} \rightarrow (7,6)
\rightarrow \textcolor{red}{(8,7)}
\end{aligned}
\]

Quindi le posizioni scelte in $v$ sono:
\[
    2 < 3 < 4 < 6 < 8
\]

e le posizioni corrispondenti in $w$ sono:
\[
    1 < 3 < 5 < 6 < 7
\]

La \textbf{sottosequenza comune} letta nelle colonne di match è:
\[
    TCTAC
\]

In questo esempio la lunghezza della longest common subsequence è quindi:
\[
    LCS(v,w)=5
\]

Anche per LCS è possibile utilizzare una \textbf{griglia di edit} di dimensione $(n+1) \times (m+1)$, con $n = |v|$ e $m = |w|$, con archi pesati e poi applicare la \textbf{programmazione dinamica}.
Tuttavia, come per la distanza di Hamming, è necessario modificare le operazioni
consentite, la funzione di costo $\delta$ e l'algoritmo di programmazione dinamica.
In particolare, dovendo trovare le posizioni dove $v$ e $w$ hanno simboli uguali,
andrà considerato il \textbf{movimento diagonale}, solo in caso di match.
Inoltre, al contrario della distanza di edit e della distanza di Hamming, serve trovare la più lunga sottosequenza comune,
quindi bisogna modificare i valori di costo per riflettere questa differenza,
escludendo i mismatch dall'allineamento.

Quindi la \textbf{funzione di costo} $\delta$ sarà definita come:
\begin{itemize}
    \item \textbf{match}: $\delta(v_i, w_j) = 1$ se $v_i = w_j$ (l'arco diagonale da $(i-1,j-1)$ a $(i,j)$)
    \item \textbf{Indel} (inserzione o delezione): $\delta(v_i,\text{-})=0$ per una delezione (l'arco verticale da $(i-1,j)$ a $(i,j)$) e $\delta(\text{-},w_j)=0$ per una inserzione (l'arco orizzontale da $(i,j-1)$ a $(i,j)$).
\end{itemize}

Per calcolare la LCS fra due sequenze si utilizza la programmazione dinamica, ma cambia l'obiettivo: invece di minimizzare il costo totale, si massimizza il numero di match, cioè la lunghezza della sottosequenza comune.

Indichiamo con $S(i,j)$ la lunghezza della LCS tra il prefisso $v[1..i]$ e il prefisso $w[1..j]$.

I \textbf{casi base} diventano:
\begin{itemize}
    \item $S(0,0)=0$
    \item $S(i,0)=0$ per $0 \leq i \leq n$
    \item $S(0,j)=0$ per $0 \leq j \leq m$
\end{itemize}

Infatti, se una delle due sequenze è vuota, non esiste alcuna sottosequenza comune non vuota, quindi la lunghezza della LCS è $0$.

Si riempie il resto della griglia con la seguente \textbf{equazione di ricorrenza}:
\begin{equation}
    S(i,j)=\max
    \begin{cases}
        S(i-1,j) & \text{(salto } v_i\text{)} \\
        S(i,j-1) & \text{(salto } w_j\text{)} \\
        S(i-1,j-1)+1 & \text{se } v_i=w_j
    \end{cases}
\end{equation}

Il valore nel nodo finale:
\begin{equation}
    S(n,m)
\end{equation}

coincide quindi con la lunghezza della longest common subsequence:

\begin{equation}
     S(n,m)=LCS(v,w)
\end{equation}

Effettuando il \textbf{backtracking} nella griglia, partendo dal nodo finale
$(n,m)$ e risalendo fino al nodo sorgente $(0,0)$, è possibile ricostruire
una matrice di allineamento $A$.

Per ottenere la \textbf{longest common subsequence}, si considerano solo le
colonne di $A$ che corrispondono a un match, cioè colonne del tipo:

\[
\begin{pmatrix}
v_i \\
w_j
\end{pmatrix}
\]

con $v_i = w_j$.

I simboli letti in queste colonne, nell'ordine dell'allineamento, formano la
longest common subsequence.

\newpage
